\chapter{Αναφορά Γραμμής Εντολών}
\label{app:cli}

Το παράρτημα αυτό αποτελεί την εξαντλητική αναφορά της διεπαφής γραμμής εντολών της υλοποίησης. Αναλυτικά παραδείγματα χρήσης με σχολιασμό δίνονται στο Κεφάλαιο~\ref{chap:usage}, ενώ η σημασιολογία των πεδίων εξόδου και των ετικετών κατάστασης τεκμηριώνεται στο Παράρτημα~\ref{app:schema}. Η διεπαφή καλείται ως \code{rubik-optimal} (εντολή κονσόλας που καταχωρίζεται κατά την εγκατάσταση) ή, ισοδύναμα, ως \code{python -m rubik_optimal.cli}. Ο Πίνακας~\ref{tab:cli-subcommands} συνοψίζει τις οκτώ υποεντολές.

\begin{table}[htbp]\centering
\small
\begin{tabular}{@{}L{0.30\linewidth}L{0.62\linewidth}@{}}
\hline
Υποεντολή & Ρόλος \\
\hline
\code{scramble} & Παραγωγή ντετερμινιστικού ανακατέματος (\eng{scramble}) από σπόρο. \\
\code{facelets} & Παραγωγή έγκυρης συμβολοσειράς ψηφίδων (\eng{facelet string})· συνώνυμο \code{random-facelets}. \\
\code{solve} & Επίλυση κατάστασης ή ανακατέματος με επιλεγόμενο επιλυτή. \\
\code{verify} & Επαλήθευση προτεινόμενης λύσης πάνω σε δοσμένη κατάσταση. \\
\code{distance} & Αναγνώριση ακριβούς απόστασης ή κάτω φράγματος. \\
\code{oracle} & Ακριβής επίλυση δέσμης καταστάσεων με κοινή φόρτωση πίνακα H48. \\
\code{benchmark} & Εκτέλεση του σώματος μετρήσεων ενός προφίλ. \\
\code{tables} & Παραγωγή πινάκων κινήσεων και πινάκων αποκοπής (\eng{pruning tables}) συντεταγμένων. \\
\hline
\end{tabular}
\caption{Οι υποεντολές της διεπαφής γραμμής εντολών \code{rubik-optimal}, όπως ορίζονται στο \code{src/rubik_optimal/cli.py}.}
\label{tab:cli-subcommands}
\end{table}

Όλες οι υποεντολές που εκτελούν αναζήτηση τυπώνουν την έξοδό τους ως JSON και ακολουθούν κοινή σύμβαση κωδικών εξόδου: η \code{solve} επιστρέφει μη μηδενικό κωδικό όταν το αποτέλεσμα φέρει ετικέτα \code{failed} ή \code{timeout}, η \code{verify} επιστρέφει μηδέν μόνο όταν η ακολουθία πράγματι λύνει την κατάσταση, και η \code{oracle} επιστρέφει μηδέν μόνο όταν όλες οι γραμμές είναι \code{exact} και επαληθευμένες. Έτσι μπορεί κανείς να βασιστεί στις εντολές μέσα σε scripts κελύφους, χωρίς να χρειάζεται να αναλύσει (parse) το JSON.

\section{Οι υποεντολές scramble και facelets}
\label{sec:cli-scramble}

\begin{verbatim}
rubik-optimal scramble [--length N] [--seed S]
rubik-optimal facelets [--length N] [--seed S]
                       [--offset K] [--json]
\end{verbatim}

Η \code{scramble} τυπώνει ντετερμινιστική ακολουθία κινήσεων μήκους \code{--length} (προεπιλογή 20) από τον σπόρο \code{--seed} (προεπιλογή 2026). Η \code{facelets} εφαρμόζει το αντίστοιχο ανακάτεμα στη λυμένη κατάσταση και τυπώνει την προκύπτουσα έγκυρη συμβολοσειρά ψηφίδων 54 χαρακτήρων. Η επιλογή \code{--offset} (προεπιλογή 0) μετατοπίζει ντετερμινιστικά την παραγωγή, ώστε ο ίδιος σπόρος να δίνει διαφορετικές καταστάσεις. Η επιλογή \code{--json} επιστρέφει επιπλέον το ανακάτεμα και τα μεταδεδομένα του (μήκος, σπόρος, μετρική, εγκυρότητα). Η υποεντολή είναι χρήσιμη όταν θέλουμε να δώσουμε απευθείας μια κατάσταση ως είσοδο, χωρίς να γνωρίζουμε την ακολουθία που την παρήγαγε.

\section{Η υποεντολή solve}
\label{sec:cli-solve}

\begin{verbatim}
rubik-optimal solve [STATE] [--solver NAME] [...]
\end{verbatim}

Το προαιρετικό όρισμα \code{STATE} δέχεται συμβολοσειρά ψηφίδων, μια ακολουθία κινήσεων που εφαρμόζεται στη λυμένη κατάσταση ή τη λέξη \code{solved} (προεπιλογή). Η είσοδος ελέγχεται αν αντιστοιχεί σε εφικτή (συναρμολογήσιμη) κατάσταση πριν κληθεί οποιοσδήποτε επιλυτής. Η επιλογή \code{--solver} δέχεται τις ακόλουθες τιμές:

\begin{itemize}
\item Πρακτικοί επιλυτές: \code{auto} (προεπιλογή· δοκιμάζει διαδοχικά την υλοποίηση δύο φάσεων σε C/C++, τον εξωτερικό προσαρμογέα και τον επιλυτή Thistlethwaite), \code{native-kociemba}, \code{thistlethwaite}, \code{korf} (η εκπαιδευτική IDA* σε Python), \code{kociemba} και \code{adapter} (ο εξωτερικός προσαρμογέας του πακέτου \code{kociemba} \cite{muodovKociembaPkg}). Στην ίδια ομάδα ανήκουν οι \code{bfs} (ρηχή εξαντλητική αναζήτηση πρώτα κατά πλάτος έως το \code{--max-depth}) και \code{inverse} (αντιστροφή της ακολουθίας εισόδου· έγκυρη μόνο όταν η είσοδος είναι ακολουθία κινήσεων).
\item Ακριβείς επιλυτές σε C/C++: \code{optimal-native} (Korf/IDA* με τον γωνιακό και τους οκτώ εξακμικούς πίνακες προτύπων), \code{h48-native} και \code{h48-oracle} (η ενσωματωμένη μηχανή H48 \cite{trontoNissyCoreH48}), \code{nissy-core-direct} (απευθείας κλήση του ενσωματωμένου πυρήνα με είσοδο κατάστασης).
\item Εξωτερικοί ακριβείς επιλυτές: \code{nissy-light} και \code{nissy-optimal} (το εξωτερικό εκτελέσιμο \code{nissy} με τους ελαφρούς και τον δημόσιο βέλτιστο πίνακά του αντίστοιχα), \code{rubikoptimal} (το εξωτερικό backend RubikOptimal).
\item Επιλυτές χαρτοφυλακίου: \code{race-optimal}, \code{resident-race-optimal} και \code{universal-optimal}, οι οποίοι συνδυάζουν τα ακριβή backends και τα πιστοποιητικά είτε παράλληλα είτε διαδοχικά (σε στάδια) (Ενότητα~\ref{sec:cli-portfolio}).
\end{itemize}

Η επιλογή \code{--max-depth} (προεπιλογή 10) ερμηνεύεται διαφορετικά από κάθε επιλυτή. Μόνο ο επιλυτής \code{korf} την τηρεί κατά γράμμα. Οι ακριβείς επιλυτές σε C/C++ την ανεβάζουν σιωπηρά σε τουλάχιστον 20. Στους επιλυτές δύο φάσεων προκύπτουν από αυτήν τα φασικά όρια (τουλάχιστον 12 και 18), ενώ ο επιλυτής Thistlethwaite ανεβάζει τα όρια των σταδίων 3 και 4 σε τουλάχιστον 13 και 15. Τιμές κάτω από αυτά τα ελάχιστα όρια ανεβαίνουν αυτόματα στο ελάχιστο, αντί να απορρίπτονται, ώστε καμία εσφαλμένη ρύθμιση να μην καθιστά κάποιον επιλυτή μη πλήρη (ατελή). Ο Πίνακας~\ref{tab:cli-solve} συνοψίζει τις υπόλοιπες γενικές επιλογές της υποεντολής.

\begin{table}[htbp]\centering
\small
\begin{tabular}{@{}L{0.34\linewidth}L{0.12\linewidth}L{0.44\linewidth}@{}}
\hline
Επιλογή & Προεπιλογή & Περιγραφή \\
\hline
\code{--timeout} & 5.0 & Όριο χρόνου αναζήτησης σε δευτερόλεπτα. \\
\code{--node-limit} & 500\,000 & Όριο κόμβων· το τηρεί μόνο ο επιλυτής \code{korf}. \\
\code{--threads} & αυτόματη & Αριθμός νημάτων για τις αναζητήσεις σε C/C++· η προεπιλογή προκύπτει από το φορτίο του μηχανήματος ή τις μεταβλητές περιβάλλοντος \code{RUBIK_OPTIMAL_THREADS} / \code{RUBIK_OPTIMAL_H48_THREADS}. \\
\code{--split-depth} & 3 & Βάθος διαμοιρασμού εργασιών στον παράλληλο IDA* σε C/C++. \\
\code{--native-child-order} & \code{heuristic-desc} & Διάταξη των κόμβων-παιδιών (επόμενων καταστάσεων) στον IDA* σε C/C++ (\code{heuristic-asc}, \code{move}). \\
\code{--dual-heuristic} & ανενεργή & Διπλή ευρετική στον Korf/IDA* σε C/C++. \\
\code{--additive-edge-pdbs} & ανενεργή & Χρήση των προσθετικών βάσεων ακμών, εφόσον έχουν παραχθεί. \\
\code{--seven-edge-pdbs} & ανενεργή & Προαιρετική χρήση των βάσεων 7 ακμών, εκτός της παγιωμένης (κλειδωμένης) διαμόρφωσης αναφοράς που χρησιμοποιεί 6 ακμές. \\
\code{--nissy-heuristic} & ανενεργή & Πρόσθετο παραδεκτό κάτω φράγμα μέσω της γέφυρας Nissy στον IDA* σε C/C++. \\
\code{--no-nissy-}\allowbreak\code{axis-transforms} & (ενεργοί) & Απενεργοποίηση των αξονικών μετασχηματισμών της γέφυρας Nissy. \\
\code{--nissy-data-dir} & N/A & Κατάλογος δεδομένων του εξωτερικού \code{nissy}. \\
\code{--upper-solution} & N/A & Υποψήφια λύση ως πιστοποιητικό άνω φράγματος για τον \code{optimal-native}. \\
\code{--upper-bound-}\allowbreak\code{proof-strategy} & \code{single-bound} & Απόδειξη βελτιστότητας με μία εξαντλητική αναζήτηση σε βάθος μικρότερο κατά ένα από το άνω φράγμα, ή με κλασικά επαναληπτικά όρια (\code{iterative}). \\
\code{--h48-start-delay} & 0.0 & Καθυστέρηση εκκίνησης του μόνιμου (\eng{resident}) H48 στους επιλυτές αγώνα, ώστε να προλαβαίνουν οι φθηνότερες απόπειρες Nissy. \\
\code{--rubikoptimal-table-dir}, \code{--rubikoptimal-}\allowbreak\code{executable}, \code{--rubikoptimal-}\allowbreak\code{package-path} & N/A & Διαδρομές πινάκων, εκτελέσιμου και πακέτου του εξωτερικού RubikOptimal. \\
\hline
\end{tabular}
\caption{Γενικές επιλογές της υποεντολής \code{solve}.}
\label{tab:cli-solve}
\end{table}

\section{Η ομάδα επιλογών H48}
\label{sec:cli-h48}

Οι επιλογές του Πίνακα~\ref{tab:cli-h48} είναι κοινές στις υποεντολές \code{solve}, \code{distance} και \code{oracle} και ρυθμίζουν τη μηχανή H48. Η σημαία \code{--h48-trusted-table} είναι ασφαλής μόνο για πίνακες που έχουν παραχθεί τοπικά και ταυτοποιηθεί από τα αποθηκευμένα μεταδεδομένα και αθροίσματα ελέγχου. Η σημαία \code{--h48-preload-table} προθερμαίνει διαδοχικά τις σελίδες του πίνακα πριν από την αναζήτηση. Προορίζεται για την πιστοποίηση δύσκολων καταστάσεων όταν ο πίνακας δεν είναι ήδη φορτωμένος στη μνήμη, και δεν επιταχύνει γενικά κάθε κλήση. Οι μετρημένες επιδράσεις των δύο σημαιών παρουσιάζονται στο Κεφάλαιο~\ref{chap:experiments}.

\begin{table}[htbp]\centering
\small
\begin{tabular}{@{}L{0.30\linewidth}L{0.14\linewidth}L{0.46\linewidth}@{}}
\hline
Επιλογή & Προεπιλογή & Περιγραφή \\
\hline
\code{--h48-solver} & ανά υποεντολή & Επίπεδο πίνακα H48 (π.χ. \code{h48h0}, \code{h48h7})· προεπιλογή \code{h48h0} στη \code{distance}, \code{h48h7} στην \code{oracle}. \\
\code{--h48-oracle} & ανενεργή & Επιλέγει το \code{h48h7}, το προφίλ βέλτιστης επίλυσης ποιότητας μαντείου (\eng{oracle}) του ενσωματωμένου πυρήνα. \\
\code{--h48-fastest} & ανενεργή & Επιλέγει τον μεγαλύτερο διαθέσιμο παραγόμενο πίνακα ποιότητας μαντείου του προφίλ. \\
\code{--h48-trusted-table} & ανενεργή & Παραλείπει την πλήρη σάρωση κατανομής του πίνακα μετά την επαλήθευση των παραγόμενων μεταδεδομένων. \\
\code{--h48-preload-table} & ανενεργή & Προθερμαίνει διαδοχικά τις σελίδες του πίνακα πριν από την αναζήτηση. \\
\code{--h48-auto-min-depth} & ανενεργή & Εκκινεί την ακριβή αναζήτηση από το παραδεκτό κάτω φράγμα H48 αντί από βάθος 0. \\
\code{--h48-profile} & \code{thesis} & Προφίλ πινάκων (\code{quick}, \code{thesis}, \code{stress}). \\
\code{--h48-table} & N/A & Ρητή διαδρομή αρχείου πίνακα H48. \\
\hline
\end{tabular}
\caption{Η κοινή ομάδα επιλογών H48 των υποεντολών \code{solve}, \code{distance} και \code{oracle}.}
\label{tab:cli-h48}
\end{table}

\section{Η υποεντολή verify}
\label{sec:cli-verify}

\begin{verbatim}
rubik-optimal verify STATE SOLUTION
\end{verbatim}

Η εντολή δέχεται την αρχική κατάσταση (συμβολοσειρά ψηφίδων ή ακολουθία) και μια υποψήφια ακολουθία λύσης, εφαρμόζει τη δεύτερη στην πρώτη και τυπώνει τα πεδία \code{ok}, \code{move_count} και \code{message}. Δεν δέχεται άλλες επιλογές.

\section{Η υποεντολή distance}
\label{sec:cli-distance}

\begin{verbatim}
rubik-optimal distance [STATE] [...]
\end{verbatim}

Η εντολή εκτελεί κλιμακωτή αναγνώριση απόστασης: εξαντλητική BFS έως το \code{--bfs-depth}, έπειτα IDA* με το συνδυασμένο παραδεκτό κάτω φράγμα των πινάκων έως το \code{--ida-depth}, και προαιρετικά τα ισχυρότερα backends. Οι ειδικές επιλογές της συνοψίζονται στον Πίνακα~\ref{tab:cli-distance}· ισχύει επιπλέον όλη η ομάδα επιλογών H48 του Πίνακα~\ref{tab:cli-h48}.

\begin{table}[htbp]\centering
\small
\begin{tabular}{@{}L{0.28\linewidth}L{0.14\linewidth}L{0.48\linewidth}@{}}
\hline
Επιλογή & Προεπιλογή & Περιγραφή \\
\hline
\code{--bfs-depth} & 5 & Βάθος της εξαντλητικής αναζήτησης πρώτα κατά πλάτος. \\
\code{--ida-depth} & 8 & Βάθος της IDA* με το συνδυασμένο κάτω φράγμα των πινάκων. \\
\code{--timeout} & 5.0 & Όριο χρόνου σε δευτερόλεπτα. \\
\code{--threads} & αυτόματη & Αριθμός νημάτων για τις αναζητήσεις σε C/C++. \\
\code{--optimal-native} & ανενεργή & IDA* σε C/C++ με τους πίνακες προτύπων μετά την BFS. \\
\code{--h48-native} & ανενεργή & Η μηχανή H48 σε C/C++ μετά την BFS, με την κατάσταση να δίνεται απευθείας ως είσοδος. \\
\hline
\end{tabular}
\caption{Ειδικές επιλογές της υποεντολής \code{distance}.}
\label{tab:cli-distance}
\end{table}

\section{Η υποεντολή oracle}
\label{sec:cli-oracle}

\begin{verbatim}
rubik-optimal oracle [STATES ...] [--input-file F] [...]
\end{verbatim}

Η εντολή υλοποιεί τη λειτουργία μαντείου: λύνει ακριβώς πολλές καταστάσεις με μία κοινή φόρτωση του πίνακα H48 και δέχεται εισόδους ως ορίσματα, από αρχείο ή από την πρότυπη είσοδο, μία ανά γραμμή. Με την επιλογή \code{--universal} χρησιμοποιείται το \code{UniversalOptimalOracle}, το οποίο δοκιμάζει πρώτα τα πιστοποιητικά μηδενικής αναζήτησης και έπειτα μία δέσμη μόνιμου H48 για τις υπόλοιπες έγκυρες εισόδους. Όταν η βελτιστότητα μιας γραμμής στηρίζεται μόνο σε ετικέτα τρίτου, η γραμμή επιστρέφεται με ετικέτα \code{external_label_exact} και ρητή αναφορά του σε τι στηρίζεται ο ισχυρισμός ακρίβειας — ποτέ ως απλό \code{exact}. Οι ειδικές επιλογές της υποεντολής συνοψίζονται στον Πίνακα~\ref{tab:cli-oracle}· ισχύουν επιπλέον οι Πίνακες \ref{tab:cli-h48} και \ref{tab:cli-portfolio}.

\begin{table}[htbp]\centering
\small
\begin{tabular}{@{}L{0.40\linewidth}L{0.11\linewidth}L{0.39\linewidth}@{}}
\hline
Επιλογή & Προεπιλογή & Περιγραφή \\
\hline
\code{STATES} (θεσιακά) & N/A & Συμβολοσειρές ψηφίδων ή ακολουθίες σε εισαγωγικά· χωρίς ορίσματα διαβάζεται η πρότυπη είσοδος. \\
\code{--input-file} & N/A & Αρχείο με μία είσοδο ανά γραμμή· κενές γραμμές και σχόλια αγνοούνται. \\
\code{--jsonl} & ανενεργή & Μία γραμμή JSON ανά είσοδο αντί ενός συγκεντρωτικού αντικειμένου. \\
\code{--stream} & ανενεργή & Διατηρεί μόνιμο backend σε C/C++ και τυπώνει κάθε αποτέλεσμα μόλις ολοκληρωθεί. \\
\code{--timeout} & 300.0 & Όριο χρόνου σε δευτερόλεπτα. \\
\code{--threads} & αυτόματη & Νήματα του backend σε C/C++. \\
\code{--max-depth} & 20 & Όριο βάθους αναζήτησης. \\
\code{--universal} & ανενεργή & Ενεργοποίηση του \code{UniversalOptimalOracle}. \\
\code{--rubikoptimal} & ανενεργή & Χρήση του εξωτερικού RubikOptimal ανά έγκυρη είσοδο, μαζί με τις αντίστοιχες επιλογές διαδρομών του. \\
\code{--resident-h48-batch-timeout} & 30.0 & Όριο ανά κατάσταση της δέσμης μόνιμου H48· αρνητική τιμή καταργεί το χωριστό όριο. \\
\code{--no-universal-}\allowbreak\code{portfolio-prepass} & ανενεργή & Παράλειψη του προπεράσματος χαρτοφυλακίου πριν από τη δέσμη μόνιμου H48. \\
\code{--universal-portfolio-}\allowbreak\code{prepass-timeout} & \code{--timeout} & Χωριστό όριο του προπεράσματος χαρτοφυλακίου/Nissy. \\
\code{--learned-certificate-log} & N/A & Προσθήκη (στο τέλος) νέων επαληθευμένων \code{exact} γραμμών σε αρχείο JSONL πιστοποιητικών. \\
\code{--no-universal-}\allowbreak\code{certificate-cache} & ανενεργή & Παράλειψη της κρυφής μνήμης πιστοποιητικών, για μετρήσεις σε πραγματική εκτέλεση (χωρίς χρήση κρυφής μνήμης). \\
\code{--no-universal-}\allowbreak\code{upper-lower-certificate} & ανενεργή & Παράλειψη της συντόμευσης πιστοποιητικού άνω/κάτω φράγματος πριν εκτελεστούν πραγματικά τα backends αναζήτησης. \\
\code{--universal-include-}\allowbreak\code{external-label-}\allowbreak\code{certificates} & ανενεργή & Επιτρέπει γραμμές των οποίων η βελτιστότητα στηρίζεται σε ετικέτα τρίτου εργαλείου· επιστρέφονται ως \code{external_label_exact}. \\
\hline
\end{tabular}
\caption{Ειδικές επιλογές της υποεντολής \code{oracle}.}
\label{tab:cli-oracle}
\end{table}

\section{Προηγμένες επιλογές χαρτοφυλακίου και συμμετριών}
\label{sec:cli-portfolio}

Οι επιλογές του Πίνακα~\ref{tab:cli-portfolio} ρυθμίζουν λεπτομερώς τους επιλυτές χαρτοφυλακίου. Καλύπτουν τα εφεδρικά περάσματα, τον παράλληλο ανταγωνισμό περιστροφών όλου του κύβου μέσα από τα διαφορετικά backends, καθώς και τις αποδείξεις με φραγμένα όρια αναζήτησης κάτω από επαληθευμένο άνω φράγμα. Είναι διαθέσιμες στην υποεντολή \code{solve} για τους επιλυτές χαρτοφυλακίου και στην υποεντολή \code{oracle} με την επιλογή \code{--universal}. Όταν η προεπιλογή είναι αρνητική ή μηδενική, η αντίστοιχη φάση παραμένει ανενεργή εκτός αν ζητηθεί ρητά.

{\small
\begin{longtable}{@{}L{0.46\linewidth}L{0.11\linewidth}L{0.34\linewidth}@{}}
\caption{Προηγμένες επιλογές χαρτοφυλακίου και αγώνων συμμετριών, κοινές στις υποεντολές \code{solve} και \code{oracle}.}
\label{tab:cli-portfolio} \\
\hline
Επιλογή & Προεπιλογή & Περιγραφή \\
\hline
\endfirsthead
\hline
Επιλογή & Προεπιλογή & Περιγραφή \\
\hline
\endhead
\hline
\multicolumn{3}{r}{\footnotesize συνέχεια στην επόμενη σελίδα} \\
\endfoot
\hline
\endlastfoot
\code{--universal-portfolio-}\allowbreak\code{fallback-timeout} & \code{--timeout} & Όριο του όψιμου εφεδρικού περάσματος χαρτοφυλακίου/Nissy. \\
\code{--universal-fallback-}\allowbreak\code{nissy-core-direct-timeout} & 10.0 & Όριο της όψιμης απευθείας εφεδρείας του ενσωματωμένου πυρήνα με είσοδο κατάστασης. \\
\code{--universal-rubikoptimal-}\allowbreak\code{fallback-timeout} & $-1$ & Εκτέλεση του RubikOptimal μετά την αποτυχία του H48 και του χαρτοφυλακίου. \\
\code{--universal-rubikoptimal-}\allowbreak\code{prepass-timeout} & $-1$ & Εκτέλεση του RubikOptimal πριν από τη δέσμη ή τον αγώνα του μόνιμου H48. \\
\code{--universal-rubikoptimal-}\allowbreak\code{race-timeout} & $-1$ & Εκτέλεση του RubikOptimal ταυτόχρονα μέσα στον αγώνα του μόνιμου H48/Nissy. \\
\code{--universal-resident-}\allowbreak\code{race-prepass-timeout} & $-1$ & Αγώνας H48/Nissy/RubikOptimal με χρονικό όριο πριν από τις διαδοχικές φάσεις. \\
\code{--universal-rubikoptimal-}\allowbreak\code{symmetry-variants} & 0 & Πλήθος περιστροφών όλου του κύβου μέσω RubikOptimal. \\
\code{--universal-rubikoptimal-}\allowbreak\code{symmetry-timeout} & \code{--timeout} & Συνολικό όριο της φάσης συμμετριών RubikOptimal. \\
\code{--universal-rubikoptimal-}\allowbreak\code{symmetry-max-concurrency} & 0 & Ταυτόχρονες διεργασίες της φάσης· με 0, οι περιστροφές λύνονται η μία μετά την άλλη μέσα σε μία ενιαία δέσμη. \\
\code{--h48-symmetry-variants} & 0 & Περιστροφές μέσω μόνιμου H48. \\
\code{--h48-symmetry-timeout} & 5.0 & Συνολικό όριο της φάσης συμμετριών H48. \\
\code{--nissy-symmetry-variants} & 0 & Περιστροφές μέσω του εξωτερικού Nissy. \\
\code{--nissy-symmetry-timeout} & \code{--timeout} & Όριο της δέσμης περιστροφών Nissy. \\
\code{--nissy-core-direct-}\allowbreak\code{symmetry-variants} & 0 & Περιστροφές μέσω του ενσωματωμένου πυρήνα με είσοδο κατάστασης. \\
\code{--nissy-core-direct-}\allowbreak\code{symmetry-timeout} & \code{--timeout} & Συνολικό όριο του αντίστοιχου αγώνα. \\
\code{--nissy-core-direct-}\allowbreak\code{symmetry-max-concurrency} & 0 & Μέγιστες ταυτόχρονες διεργασίες· με 0 εκκινούν όλες οι περιστροφές. \\
\code{--h48-parallel-symmetry-variants} & 0 & Παράλληλος αγώνας περιστροφών μέσω H48. \\
\code{--h48-parallel-symmetry-timeout} & 5.0 & Συνολικό όριο του παράλληλου αγώνα H48. \\
\code{--h48-parallel-symmetry-}\allowbreak\code{max-concurrency} & 0 & Μέγιστες ταυτόχρονες διεργασίες H48. \\
\code{--h48-parallel-symmetry-}\allowbreak\code{order-by-lower-bound} & ανενεργή & Διάταξη των υποψήφιων περιστροφών με φθηνή δέσμη παραδεκτών κάτω φραγμάτων H48. \\
\code{--h48-parallel-symmetry-}\allowbreak\code{lower-bound-order-timeout} & 30.0 & Όριο της δέσμης κάτω φραγμάτων για τη διάταξη. \\
\code{--h48-lower-bound-}\allowbreak\code{symmetry-variants} & 23 & Πλήθος περιστροφών στις οποίες υπολογίζονται κάτω φράγματα H48 πριν γίνουν δεκτά τα πιστοποιητικά άνω/κάτω φράγματος. \\
\code{--kociemba-upper-bound-}\allowbreak\code{symmetry-variants} & 23 & Περιστροφές μέσω Kociemba ως βελτιωμένα άνω φράγματα πριν από την αποδοχή πιστοποιητικών. \\
\code{--h48-upper-bound-proof-timeout} & 0 & Χρονικό όριο (σε δευτερόλεπτα) για την απόδειξη μέσω H48 ότι δεν υπάρχει λύση κάτω από το επαληθευμένο άνω φράγμα. \\
\code{--h48-upper-bound-proof-max-gap} & 1 & Μέγιστη διαφορά άνω και κάτω φράγματος για να εκτελεστεί η απόδειξη. \\
\code{--native-korf-}\allowbreak\code{upper-bound-proof-timeout} & 0 & Χρονικό όριο (σε δευτερόλεπτα) για την απόδειξη μέσω Korf/IDA* σε C/C++ ότι δεν υπάρχει λύση ένα βάθος κάτω από την επαληθευμένη λύση. \\
\code{--native-korf-}\allowbreak\code{upper-bound-proof-max-gap} & 1 & Αντίστοιχη μέγιστη διαφορά (άνω–κάτω φράγματος) για να εκτελεστεί η απόδειξη σε C/C++. \\
\end{longtable}
}

\section{Οι υποεντολές benchmark και tables}
\label{sec:cli-benchmark}

\begin{verbatim}
rubik-optimal benchmark [--profile P] [--seed S]
                        [--quick] [--root DIR]
rubik-optimal tables    [--profile P] [--seed S]
                        [--quick] [--root DIR]
\end{verbatim}

Η \code{benchmark} εκτελεί το σώμα μετρήσεων του προφίλ \code{--profile} (τιμές \code{quick}, \code{thesis}, \code{stress}) με σπόρο \code{--seed} (προεπιλογή 2026) και τυπώνει τις διαδρομές των παραγόμενων αρχείων· η σημαία \code{--quick} είναι συντόμευση για το προφίλ \code{quick}. Η \code{tables} παράγει τους πίνακες κινήσεων και αποκοπής των συντεταγμένων με τις ίδιες επιλογές. Η επιλογή \code{--root} ορίζει τον κατάλογο εργασίας (προεπιλογή ο τρέχων).

\section{Σενάρια παραγωγής τεκμηρίων}
\label{sec:cli-scripts}

Η διεπαφή γραμμής εντολών εξυπηρετεί άμεσους ελέγχους και χειροκίνητη διερεύνηση, ενώ τα σενάρια του καταλόγου \code{scripts/} παράγουν τα τεκμήρια αναφοράς με σταθερά ονόματα αρχείων και συμβάσεις προφίλ. Οι αριθμοί του Κεφαλαίου~\ref{chap:experiments} προέρχονται αποκλειστικά από τα σενάρια αυτά, με την πλήρη αλληλουχία εκτέλεσης του Παραρτήματος~\ref{app:repro}.

Το \path{scripts/generate_tables.py} παράγει τους πίνακες συντεταγμένων και πρέπει να εκτελείται όταν αλλάζουν οι κωδικοποιήσεις συντεταγμένων, οι συμβάσεις κινήσεων ή οι προδιαγραφές των πινάκων. Τα \path{scripts/generate_corner_pdb.py} και \path{scripts/generate_edge_pdb.py} παράγουν τον γωνιακό πίνακα προτύπων (\eng{pattern database}, PDB) και τους πίνακες ακμών αντίστοιχα. Το \path{scripts/generate_h48_tables.py} παράγει τους πίνακες H48 με τα μεταδεδομένα και τα αθροίσματα ελέγχου τους: η επιλογή \code{--solver h48h0} παράγει τον μικρό πίνακα, ενώ η επιλογή \code{--oracle} τον πίνακα \code{h48h7} που χρησιμοποιείται στις γραμμές πίεσης ποιότητας μαντείου (για το ότι το διατηρημένο τεκμήριο δεν αναπαράγεται πανομοιότυπα, byte προς byte, βλ.\ Ενότητα~\ref{sec:repro-h48h7}).

Το \code{scripts/run_benchmarks.py} εκτελεί το σώμα μετρήσεων του κύβου 3×3×3 και πρέπει να εκτελείται όταν αλλάζουν επιλυτές, ευρετικές, όρια χρόνου ή περιπτώσεις. Το \code{scripts/run_3x3_optimal.py} εκτελεί τους ακριβείς επιλυτές, και με τις επιλογές \code{--backend h48-native} και \code{--h48-oracle} χρησιμοποιεί τον πίνακα \code{h48h7}. Το \code{scripts/generate_pocket_cube.py} παράγει την πλήρη κατανομή του κύβου 2×2×2 και τις αντιπροσωπευτικές γραμμές του. Εξειδικευμένα σενάρια καλύπτουν τις επιμέρους μετρήσεις του Κεφαλαίου~\ref{chap:experiments}: πιστοποιούν τη μηχανή μαντείου σε ένα μικρό σώμα μετρήσεων που περιλαμβάνει το superflip, μετρούν την επιτάχυνση από την κοινή φόρτωση του πίνακα, απομονώνουν το κόστος του πλήρους ελέγχου του πίνακα και καταγράφουν τη δημόσια υποεντολή \code{oracle} της γραμμής εντολών.

Το \code{scripts/generate_figures.py} ανανεώνει τους πίνακες LaTeX και τα αρχεία δεδομένων σχημάτων που εξαρτώνται από τις συνόψεις των μετρήσεων· το όνομά του κρατήθηκε από παλιότερες εκδόσεις, αν και πλέον παράγει και πίνακες (όχι μόνο σχήματα). Το \code{scripts/verify_results.py} εκτελείται μετά από κάθε παραγωγή μετρήσεων ή τεκμηρίων του 2×2×2, όπως περιγράφεται στο Παράρτημα~\ref{app:repro}.

\section{Ερμηνεία της εξόδου}
\label{sec:cli-output}

Η έξοδος κάθε εντολής επίλυσης κρατά χωριστά την ετικέτα κατάστασης και την ένδειξη επαλήθευσης· τα δύο πεδία πρέπει να διαβάζονται μαζί. Το \code{verified=true} σημαίνει ότι η λύση λύνει πράγματι τον κύβο, ενώ μόνο η ετικέτα \code{exact} δηλώνει αποδεδειγμένα βέλτιστη λύση. Μια ετικέτα \code{non_exact} συνοδεύει έγκυρη λύση χωρίς απόδειξη ελάχιστου μήκους, ενώ η ετικέτα \code{timeout} δηλώνει ότι δεν επιστράφηκε λύση εντός του ορίου· το πεδίο \code{notes} εξηγεί τότε την αιτία. Ο πλήρης κατάλογος των ετικετών και η ακριβής σημασιολογία τους δίνονται στο Παράρτημα~\ref{app:schema}.
